
\documentclass[UTF8,zihao=5,a4paper]{ctexart}
\usepackage[a4paper,top=2.5cm,bottom=2.5cm,left=2.5cm,right=2.5cm]{geometry}
\usepackage{amsmath,amssymb,bm}
\usepackage{graphicx}
\usepackage{booktabs}
\usepackage{caption}
\usepackage{enumitem}
\usepackage{listings}
\usepackage{xcolor}
\usepackage{setspace}
\usepackage{array}
\usepackage{tabularx}
\usepackage{fancyhdr}
\usepackage{xurl}

\setCJKmainfont{Noto Serif CJK SC}
\setCJKsansfont{Noto Sans CJK SC}
\setCJKmonofont{Noto Sans Mono CJK SC}
\setmainfont{Liberation Serif}
\setmonofont{Noto Sans Mono CJK SC}[Scale=0.82]

\onehalfspacing
\setlength{\parindent}{2em}
\captionsetup{skip=4pt}
\emergencystretch=3em
\sloppy

\pagestyle{fancy}
\fancyhf{}
\renewcommand{\headrulewidth}{0pt}
\fancyfoot[C]{\thepage}

\ctexset{
  section/format = {\centering\zihao{4}\bfseries},
  subsection/format = {\raggedright\zihao{5}\bfseries},
  section/beforeskip = 12pt plus 2pt minus 2pt,
  section/afterskip = 8pt,
}

\lstdefinestyle{pseudo}{
  basicstyle=\ttfamily\zihao{-5},
  breaklines=true, breakatwhitespace=false,
  columns=fullflexible, keepspaces=true,
  frame=single, rulecolor=\color{gray!55},
  framesep=4pt, xleftmargin=6pt, xrightmargin=2pt,
  showstringspaces=false, extendedchars=true,
}
\lstdefinestyle{pycode}{
  language=Python,
  basicstyle=\ttfamily\zihao{6},
  keywordstyle=\color{blue!65!black},
  commentstyle=\color{green!42!black}\itshape,
  stringstyle=\color{orange!75!black},
  breaklines=true, breakatwhitespace=false,
  breakindent=12pt, postbreak=\mbox{\textcolor{gray}{$\hookrightarrow$}\space},
  columns=fullflexible, keepspaces=true,
  showstringspaces=false, extendedchars=true,
  numbers=left, numberstyle=\tiny\color{gray}, numbersep=6pt,
  frame=leftline, rulecolor=\color{gray!45}, framesep=5pt,
  xleftmargin=17pt, tabsize=4, upquote=true,
}

% 手工编号的图表题注（与正文交叉引用一致）
\captionsetup[figure]{labelformat=empty,justification=centering}
\captionsetup[table]{labelformat=empty,justification=centering}

\begin{document}

\begin{center}
{\zihao{3}\bfseries 水泥烧成系统电除尘器的协同优化控制}
\end{center}
\vspace{6pt}
\begin{center}{\zihao{4}\bfseries 摘\quad 要}\end{center}
\vspace{2pt}
本文先分辨附件中哪些关系能由数据确定、哪些不能，再据此回答四问。

\textbf{能确定的是总功率}：$P=73.1+0.0850\sum_i U_i^{2}+5.39\times10^{4}\sum_i T_i^{-1}$，在未参与拟合的测试日上$R^2=0.993$、RMSE=5.98 kW、平均偏差−0.02 kW。第一项是电场功率，与电压平方成正比；第二项是振打功率，与振打频率$60/T_i$成正比。由此得到一个易被忽略的事实：\textbf{振打环节占总功率的38.8\%}，与电压同属一级节电杠杆；入口温度、浓度、流量则几乎不额外影响总功率。

\textbf{问题一}分两层回答。能确定的部分：上述功率律、历史操作规律（入口浓度升高时四个电场的振打周期一律缩短，$r=-0.21\sim-0.28$，四个时段方向一致）以及两段持续异常工况的控制响应。不能确定的部分是出口浓度——10030条有效记录全部落在48.74～50.00 mg/Nm³，其中5491条恰为50，与全部12个变量的相关系数绝对值都小于0.017，属传感器量程截顶，故排放层只能按Deutsch–Anderson机理推演并逐处标注。振打峰值不再假定指数：由“每个周期积下的灰量正比于$L\cdot T$”可推出单次峰值随周期近似线性放大，而它对小时均值的贡献$\propto L\cdot T\cdot(1/T)=L$，\textbf{与周期无关}。

\textbf{问题二}：按入口三变量把运行状态分为4类（轮廓系数0.4084）。以记录值50 mg/Nm³为起点、不作任何缩放，在历史运行范围内求解，四类工况达到10 mg/Nm³的最低功率为1851.1～2104.3 kW，加权1973.0 kW，比历史实际的1783.2 kW\textbf{高10.6\%}（日电耗42.8$\rightarrow$47.4 MWh）。收紧排放必然多耗电，不存在“排放不变还省电”。

\textbf{问题三}：每多花一份电功率能换来多少除尘效果，可写成$\alpha_i/(b_iU_{i}^{ref\,2})$。第三电场的参考电压最低，性价比在四类工况中均居首；原先“前级优先”的说法其实由排序约束$U_1\ge U_3+3$带出，而非权重$\alpha_i$本身。

\textbf{问题四}：闭式解$\Delta P=P_{field}\ln[(10-B)/(5-B)]/K_{10}$给出123.1～142.4 kW，与数值解相差不足1\%；加权功率由1973.0升至2105.5 kW，\textbf{增幅6.72\%}，日电耗多3.2 MWh。增幅与$2/p$成正比（$K\propto U^p$），$p$是全文唯一需要现场标定的数，一次电压阶跃试验即可测出。更关键的限制不是电耗而是\textbf{电源容量}：10 mg/Nm³仍在历史电压范围内，5 mg/Nm³则需各场电压比历史最高值再高最多5\%；若要求每一步都落在历史做过的操作组合里，八组解中七组仍可行，代价是多耗1.5\%～6.1\%电力。

\vspace{8pt}

\noindent \textbf{关键词：} 电除尘器；可识别性审计；Deutsch–Anderson灰箱；振打再飞扬；约束优化
\clearpage

\section{问题重述与分析}

\subsection*{1.1 问题背景}

电除尘器通过高压电场使粉尘荷电并向收尘极迁移。提高电场电压通常能够增强荷电与迁移，但会提高运行功率并受到火花、绝缘和联锁条件限制；振打过于频繁会增加再飞扬次数，周期过长又可能导致积灰加重和单次脱落量增大。水泥工业颗粒物排放受到国家标准约束[1]，题目进一步给出10 mg/Nm³和5 mg/Nm³两个控制目标，要求同时考虑排放、总功率和振打峰值。

本文研究对象的四电场电除尘器结构及其可测量与可调量如图1所示：含尘烟气依次流经四个电场，每个电场由独立高压电源施加电压$U_i$并由独立振打装置按周期$T_i$清灰，入口侧可测入口温度$H$、入口粉尘浓度$C_{in}$与烟气流量$Q$，出口侧可测出口粉尘浓度$C_{out}$与总功率$P$。

\begin{figure}[htbp]
\centering
\includegraphics[width=0.96\textwidth]{figures_paper/00_esp_schematic.pdf}
\caption*{\normalsize 图1　四电场电除尘器结构与变量定义}
\end{figure}

题目包含四项相互依赖的任务：第一，解释入口温度、入口浓度、烟气流量、四电场电压和振打周期与出口浓度之间的关系，并分析振打峰值；第二，划分典型工况，在10 mg/Nm³约束下寻找最低功率控制参数；第三，选取差异显著的工况解释控制优先级；第四，将约束收紧至5 mg/Nm³并估计功率代价。

\subsection*{1.2 总体思路}

附件里有两列结果变量，信息量却相差悬殊：总功率几乎完全由控制量决定，而出口浓度被压在量程上限附近、几乎不动。如果不加区分地把两者都当成建模目标，模型在出口浓度上学到的只会是传感器的封顶读数。本文因此先分层、后优化：

\begin{enumerate}[itemsep=1pt,topsep=3pt]
\item 逐列检查哪些关系能被数据辨认出来，把“可由附件验证的”与“只能靠机理推演的”分开；
\item 对可识别的总功率，按物理形式（电场功率$\propto$$U^2$、振打功率$\propto$$1/T$）建立并检验模型；
\item 对不可识别的排放层，以\textbf{记录值本身}为锚点建立Deutsch–Anderson灰箱，不作任何未经证实的量纲改写；
\item 在历史运行范围内求解优化问题，并把“结果在哪些方面超出了历史经验”当作结论的一部分如实报告。
\end{enumerate}

第3点是本文与常见做法的关键差别。附件记录的出口浓度约50 mg/Nm³，题目要求的却是10与5 mg/Nm³，相差5～10倍。承认这个差距，题目就自然变成“从现状降到超低排放要怎么调、要多花多少电”。反过来，若为了让数据靠近目标区间而假设记录值点错了小数点，第二问就会变成“把排放放宽到10 mg/Nm³能省多少电”——与题意正好相反。

\section{数据审计：哪些关系能从附件中辨认出来}

\subsection*{2.1 数据完整性与派生变量}

附件覆盖2024年5月1日0时至5月7日23时59分，共10080条记录，采样间隔1 min，时间戳无重复、无缺口。原始字段包括入口温度$H_t$、入口浓度$C_{in,t}$、烟气流量$Q_t$、四电场电压$U_{i,t}$、四电场振打周期$T_{i,t}$、出口浓度$C_{out,t}$和总功率$P_t$。出口浓度缺失50条（0.50\%），其余字段完整。

为表达过程负荷和振打风险，构造粉尘质量负荷

\begin{equation}
L_t=C_{in,t}Q_t/1000
\end{equation}

其中$C_{in}$单位为g/Nm³，$Q$单位为Nm³/h，故$L$单位为kg/h。另构造电压平方和、振打频率$60/T_i$与每周期积尘代理$L_tT_i/3600$。所有聚类与回归参数仅在训练期（5月1—5日）估计，验证期（5月6日）与回顾性测试期（5月7日）不参与拟合，划分按时间顺序而非随机打乱。

\subsection*{2.2 逐列审计结论}

对14列逐一检查的结果见表1：\textbf{只有出口浓度一列不能用，其余各列信息量都够}。


\begin{table}[htbp]
\centering
\caption*{\normalsize 表1　附件逐列可识别性审计结论}
\vspace{-2pt}
\zihao{-5}
\begin{tabularx}{\textwidth}{l>{\hsize=1.466\hsize\raggedright\arraybackslash}X>{\hsize=0.534\hsize\raggedright\arraybackslash}X}
\toprule
数据块 & 审计证据 & 可用性 \\
\midrule
$C_{out}$ 出口浓度 & 有效值10030条全部位于48.74～50.00 mg/Nm³，5491条恰等于50；与全部12个变量的同期相关系数绝对值均$<0.017$，1～120 min滞后同样无稳定关系；封顶分类AUC=0.5098，未封顶子样本回归验证$R^2=-0.019$ & \textbf{不可识别，不可修复} \\
$P$ 总功率 & 由控制量近乎确定，残差标准差约6 kW（见§5.2） & 精确可辨识 \\
$H,C_{in},Q$ 工况变量 & 连续波动并含两段持续异常事件（§4.3） & 工况分层与扰动分析 \\
$U_i,T_i$ 控制设定 & 与工况变量存在稳定相关结构（§4.2） & 可提取历史操作规律 \\
\bottomrule
\end{tabularx}
\end{table}

出口浓度的分布与截顶情况见图2。\textbf{这里必须说清楚一点}：以上结果只说明“这份记录里没有可供学习的信号”，不能反过来说物理变量对真实排放没有影响。已有研究曾用宽负荷下的机组负荷与各电场二次电流训练出口浓度模型[8]；相比之下，本附件既没有二次电流，出口记录本身也几乎不变化，不具备同样的条件。附录B给出直接以$C_{out}$为目标建模的对照结果，可与其他方案交叉比较。

\begin{figure}[htbp]
\centering
\includegraphics[width=0.92\textwidth]{figures_paper/01_outlet_data_diagnosis.png}
\caption*{\normalsize 图2　出口浓度数据诊断}
\end{figure}

首24 h过程量的变化如图3所示。入口变量和功率具有连续变化和工况迁移特征，而出口浓度记录几乎被压缩在单点附近。

\begin{figure}[htbp]
\centering
\includegraphics[width=0.92\textwidth]{figures_paper/02_process_timeseries.png}
\caption*{\normalsize 图3　首24 h过程量时序}
\end{figure}

\subsection*{2.3 排放锚点的取法}

本文\textbf{不对出口记录作任何缩放}，直接把各工况出口浓度的记录中位数（都是50 mg/Nm³）当作机理模型的起算点（下称锚点）。这样一来，题目的定位就很清楚：设备目前离超低排放还很远，本文要解的是“从现状降到10与5 mg/Nm³该怎么调、代价多大”。这个取法带来三个后果，后文都会明说：最优点必然落在历史电压范围的上端（§5.5）；10$\rightarrow$5 mg/Nm³的功率增幅明显小于从更低起点出发时的估计（§7）；排放层的数值仍属机理推演，须在补采未截顶的排放数据后重新标定（附录C）。

\section{模型假设、符号与总体框架}

\subsection*{3.1 模型假设}

本文将假设分为数据事实、情景假设、优化边界与运行约束四类，并逐条标注其证据状态，如表2所示。


\begin{table}[htbp]
\centering
\caption*{\normalsize 表2　模型假设及证据状态}
\vspace{-2pt}
\zihao{-5}
\begin{tabularx}{\textwidth}{l>{\hsize=1.127\hsize\raggedright\arraybackslash}X>{\hsize=0.873\hsize\raggedright\arraybackslash}X}
\toprule
类别 & 假设 & 证据状态 \\
\midrule
数据事实 & 7天记录按分钟连续，除50条出口浓度缺失外不填造目标值 & 附件可核验 \\
数据事实 & 总功率按$U_i^2$与$1/T_i$的物理形式建模，按时间滚动验证 & 可由验证与测试结果检查 \\
数据事实 & 排放锚点取出口浓度记录中位数50 mg/Nm³，不作缩放 & 附件记录值 \\
情景假设 & 去除指数按$K=K_0\,S(U)$随电压幂次$p$缩放，中心取$p=2$ & Deutsch迁移速度形式，未经现场标定 \\
情景假设 & 电场贡献权重$\alpha_i$取等权、弱前级与中心前级三种形状 & 工程先验与消融 \\
情景假设 & 振打峰值指数$\gamma=1$由积灰量正比关系推导，扫描1.00/1.30/1.60 & 本文推导＋结构情景 \\
优化边界 & 控制量位于各工况训练期实测最小—最大范围内 & 历史运行包络，不等同设备安全边界 \\
优化边界 & 5 mg/Nm³档允许电压上界外推至历史最大值的1.05倍 & 显式声明的外推，须经容量核算 \\
运行约束 & 前级电压高于后级至少3 kV，后级振打周期长于前级至少100 s & 工程排序约束 \\
\bottomrule
\end{tabularx}
\end{table}

\subsection*{3.2 主要符号}


\begin{table}[htbp]
\centering
\caption*{\normalsize 表3　主要符号与单位}
\vspace{-2pt}
\zihao{-5}
\begin{tabular}{lll}
\toprule
符号 & 含义 & 单位 \\
\midrule
$U_i$ & 第$i$电场电压 & kV \\
$T_i$ & 第$i$电场振打周期 & s \\
$C_{in},C_{out}$ & 入口、出口粉尘浓度 & g/Nm³、mg/Nm³ \\
$Q$ & 烟气流量 & Nm³/h \\
$L$ & 粉尘质量负荷 & kg/h \\
$P,\widehat P$ & 实际、预测总功率 & kW \\
$b_i,c_i$ & 电场功率系数、振打功率系数 & kW/kV²、kW·s \\
$K,K_0$ & 去除指数及其历史锚点值 & — \\
$S(U)$ & 电压强度因子，$U=U^{ref}$时为1 & — \\
$p$ & 去除指数对电压的幂次，$K\propto U^p$ & — \\
$B(T)$ & 振打引起的峰值增量代理 & mg/Nm³ \\
$\rho,\gamma$ & 振打相位叠加尺度、峰值—周期指数 & — \\
\bottomrule
\end{tabular}
\end{table}

框架上，工况与总功率由实测数据辨识，排放约束由显式先验建立，两者仅在支持域优化阶段合并；本文对每一个数值结论都标注其属于“实测辨识”还是“条件式推演”。

\section{问题一：可识别的影响关系与振打峰值}

第一问要求解释各变量与出口浓度的关系，并分析振打峰值。按§2的审计结论，本节分两层作答：先给出附件\textbf{能够}支持的三类实证结论，再给出只能靠机理推演的排放部分。

\subsection*{4.1 总功率的生成律（实证）}

在训练期以物理形式拟合，得到

\begin{equation}
P=73.09+\sum_{i=1}^{4}b_iU_i^{2}+\sum_{i=1}^{4}\frac{c_i}{T_i}
\end{equation}

其中$b=(0.0849,0.0851,0.0847,0.0853)$ kW/kV²、$c=(5.43,5.37,5.49,5.26)\times10^{4}$ kW·s。四个电场的$b_i$彼此相差不到0.8\%，说明四台高压电源的电气特性一致；写成合并形式即$P\approx73.0+0.0850\sum U_i^2+5.39\times10^{4}\sum T_i^{-1}$。这个式子的物理含义很直白：\textbf{电场功率与电压平方成正比，振打功率与振打频率$60/T_i$成正比}。模型的检验结果见§5.2。

从这个式子还能直接读出一个此前被忽略的事实。在训练期的中位设定下，电场功率为1020.3 kW，振打功率为693.7 kW，即\textbf{振打环节占总功率的38.8\%}。振打周期在题面里看似配角，对电耗的影响却和电压一个量级。相应地，把四个电场的周期从中位延长到历史P95可省下56.1 kW，延长到历史最大值可省下109.1 kW，代价是单次积灰量分别增加8.4\%和17.9\%（图5b）。§7.3会说明，这笔可省的电几乎正好抵掉排放限值收紧一档所增加的电耗——这正是“不能为省电而无限延长振打周期”的定量依据。

还有一点值得指出：把入口温度、浓度、流量也加进这个式子后，验证集与测试集的$R^2$变化都不到$10^{-3}$（见表8末行）。\textbf{总功率只由控制量决定，与工况无关}。后文按工况分别优化得到的功率差异，全部来自各工况不同的排放起算点和取值范围，而不是工况本身直接影响了功率。

\subsection*{4.2 历史操作规律（实证）}

出口浓度虽然辨认不出来，但这不等于附件里没有关于“操作策略”的信息。图4(c)给出整周记录中工况变量与控制设定的相关系数，其中有两类规律值得报告：

\textbf{规律一（稳定）}：入口浓度升高时，四个电场的振打周期一律缩短（$r=-0.208$、$-0.211$、$-0.277$、$-0.278$）。该方向在训练期、验证期与测试期三个时窗内符号一致，是附件中可复现的操作规律，与“高负荷下不宜延长周期”的工程直觉相符。

\textbf{规律二（不稳定）}：入口浓度升高时，前两个电场的电压随之升高（整周$r=+0.172,+0.166$），后两个电场几乎不动（$r=+0.019$）。但这条规律并不贯穿整周：$r(C_{in},U_1)$在前5天是$-0.052$，到第6、7天分别变成$+0.507$和$+0.492$。也就是说，“浓度高就抬前级电压”是最后两天才出现的做法。本文如实报告这一点：前后级响应的不对称在整周和最后两天都成立，却不能当作贯穿全周的稳定规律。这也说明历史数据中至少发生过一次操作策略调整，因此第三问只把它当作参照，不当作定论。


\begin{table}[htbp]
\centering
\caption*{\normalsize 表4　工况变量与控制设定的历史相关系数（全记录期）}
\vspace{-2pt}
\zihao{-5}
\begin{tabular}{lrrr}
\toprule
控制量 & 与入口温度$H$ & 与入口浓度$C_{in}$ & 与烟气流量$Q$ \\
\midrule
$U_1$ & −0.130 & +0.172 & +0.074 \\
$U_2$ & −0.127 & +0.166 & +0.070 \\
$U_3$ & −0.112 & +0.019 & −0.068 \\
$U_4$ & −0.113 & +0.019 & −0.063 \\
$T_1$ & −0.075 & −0.208 & −0.180 \\
$T_2$ & −0.064 & −0.211 & −0.175 \\
$T_3$ & −0.159 & −0.277 & −0.130 \\
$T_4$ & −0.151 & −0.278 & −0.135 \\
\bottomrule
\end{tabular}
\end{table}

\subsection*{4.3 两段持续异常工况（实证）}

附件含两段持续数十分钟以上的异常工况，其控制端响应方向相反，构成理解设备行为的两个直接入口（图4a、4b）。


\begin{table}[htbp]
\centering
\caption*{\normalsize 表5　两段持续异常工况及控制端响应}
\vspace{-2pt}
\zihao{-5}
\begin{tabularx}{\textwidth}{l>{\hsize=0.824\hsize\raggedright\arraybackslash}X>{\hsize=1.061\hsize\raggedright\arraybackslash}X>{\hsize=1.115\hsize\raggedright\arraybackslash}X}
\toprule
事件 & 时间窗 & 特征 & 控制端响应 \\
\midrule
入口浓度尖峰 & 05-02 14:02—15:49，其中41 min超过60 g/Nm³ & 入口浓度峰值72.47 g/Nm³（基线中位37.41） & $U_1$均值63.7 kV（峰值67.8），振打周期同步缩短，总功率均值1924 kW、峰值2020 kW \\
入口温度尖峰 & 05-06 13:00—14:59，持续120 min & 入口温度峰值158.2 $^\circ$C（基线中位126.5），同期入口浓度降至均值25.9 g/Nm³ & $U_1$均值降至55.5 kV，总功率均值降至1668 kW \\
\bottomrule
\end{tabularx}
\end{table}

第二段事件在工程上更值得注意：\textbf{140～160 $^\circ$C正是水泥窑尾粉尘比电阻最高的温度区间}。比电阻过高会引发反电晕，使粉尘向极板迁移的速度下降、火花率上升，是电除尘最难运行的一段工况[3-4]。附件在这段时间里的记录显示，操作端选择了降压而不是升压，这与“高比电阻区应限制电流密度以抑制反电晕”的常规处置一致；但附件没有二次电流和火花率记录，无法判断这是否为最优处置。由此可得两点：第一，第二问的工况划分把温度作为主轴之一是必要的；第二，本文四类工况的温度中心约为120.7与130.4 $^\circ$C，\textbf{都没有覆盖这段高比电阻区间}，因此本文结论不能外推到该区间，须单独补采数据重新标定（§9.3）。

\begin{figure}[htbp]
\centering
\includegraphics[width=0.92\textwidth]{figures_paper/04b_operating_evidence.png}
\caption*{\normalsize 图4　附件中可识别的运行证据}
\end{figure}

\subsection*{4.4 排放响应的灰箱骨架（条件式推演）}

Deutsch–Anderson关系用指数形式描述电除尘效率。它的理想化形式忽略了再飞扬、气流分布和粒径变化，因此只能当作\textbf{骨架}：形式取自机理，系数则须另行给定（本文称“灰箱”）[2-4]。White的经典专著[12]与Parker的电除尘器电气运行专著[13]都强调，电气运行、粉尘性质与振打三者是相互耦合的。定义去除指数

\begin{equation}
K=\ln\frac{1000C_{in}}{C_{out}},\qquad C_{out}=1000C_{in}e^{-K}
\end{equation}

$K$越大表示除尘越彻底：$K$每增加$\ln 2$，出口浓度就减半。对工况$k$，以历史中位电压$U_{i,k}^{ref}$和出口浓度记录中位数$C_{0,k}=50$ mg/Nm³为起算点，得$K_{0,k}=\ln(1000C_{in,k}/C_{0,k})$，四类工况依次为6.257、6.514、6.702、6.797。粉尘向极板迁移的速度与电场强度的平方近似成正比，故取

\begin{equation}
K_k(U,T)=K_{0,k}\,S_k(U)-\{G_k(T)-G_k(T^{ref})\},\qquad
S_k(U)=\frac{\sum_{i}\alpha_i\left(U_i/U_{i,k}^{ref}\right)^{p}}{\sum_i\alpha_i}
\end{equation}

\begin{equation}
G_k(T)=\sum_{i=1}^{4}\beta_{i,k}
\left(\frac{T_i}{T_{i,k}^*}+\frac{T_{i,k}^*}{T_i}-2\right),
\qquad \beta_{i,k}=\beta_i\left(\frac{L_k}{L_{med}}\right)^{1/2}
\end{equation}

这样构造有一个好处：当$U=U^{ref}$、$T=T^{ref}$时式子严格还原为$K=K_{0,k}$，即\textbf{模型天然复现历史运行点}，因此不再需要一个额外的整体比例参数$\eta$，全部电压敏感性都归结到一个幂次$p$上（中心取$p=2$）。$G_k(T)$在$T_i=T_{i,k}^*$处取最小值，表示周期过短或过长都不利：过短会频繁引起再飞扬，过长则积灰加重、单次脱落量变大。试验研究也确实观察到再飞扬与最大允许周期之间存在这种权衡[6]。再加一点安全裕度$s$，连续排放的基础值为$C_{base,k}=1000C_{in,k}\exp[-K_k(U,T)+s]$。


\begin{quote}\small 本节全部排放系数均为工程先验，不是由附件出口浓度估计得到的参数。\end{quote}

\begin{table}[htbp]
\centering
\caption*{\normalsize 表6　灰箱排放模型参数来源与证据状态}
\vspace{-2pt}
\zihao{-5}
\begin{tabularx}{\textwidth}{lr>{\hsize=1.013\hsize\raggedright\arraybackslash}X>{\hsize=0.987\hsize\raggedright\arraybackslash}Xc}
\toprule
参数 & 中心取值 & 含义 & 来源 & 是否由附件识别 \\
\midrule
$p$ & 2.0 & 去除指数对电压的幂次 & Deutsch迁移速度形式 & 否 \\
$\alpha_i$ & (1.15,1.10,0.95,0.80) & 电场去除贡献权重 & 工程先验与消融 & 否 \\
$\beta_i$ & (0.30,0.27,0.20,0.17) & 周期偏离损失 & 工程情景 & 否 \\
$b_0$ & 1.2 & 峰值尺度 & 工程情景 & 否 \\
$\gamma$ & 1.0 & 周期—峰值指数 & 本文由积灰量推导 & 否 \\
$\rho$ & 1.0 & 相位叠加尺度 & 同相保守情景 & 否 \\
$s$ & 0.08 & 安全裕度 & 工程情景 & 否 \\
\bottomrule
\end{tabularx}
\end{table}

\subsection*{4.5 振打峰值：由积灰量推导而非假定}

附件只给出振打周期的设定值，没有实际振打时刻、相位和秒级浓度，所以逐次峰值本身无法辨认。但峰值\textbf{随周期怎样变化}可以由质量守恒直接推出，不必假定指数：设单位时间的积灰速率正比于粉尘负荷$L$，那么一个周期内积下的灰饼质量就正比于$L\cdot T$；一次振打抛落的粉尘量又正比于灰饼质量，于是

\begin{equation}
B_k(T;\rho,\gamma)=b_0\rho\left(\frac{L_k}{L_{med}}\right)
\sum_{i=1}^{4}w_i\left(\frac{T_i}{T_{i,k}^*}\right)^{\gamma},\qquad \gamma=1 
\end{equation}

由此得到两条互补的结论，也正是第一问关于振打峰值的完整回答：

\begin{itemize}[itemsep=1pt,topsep=3pt]
\item \textbf{单次瞬时峰值随周期近似线性放大}：周期延长10\%与20\%时，单次峰值代理相应放大10\%与20\%；
\item \textbf{对小时均值的贡献与周期无关}：单位时间的总再飞扬量正比于$L\cdot T\times(1/T)=L$，只取决于粉尘负荷。
\end{itemize}

换句话说，延长振打周期并不改变再飞扬对小时均值的贡献，只是把同样多的粉尘集中到次数更少、单次更高的尖峰里释放。这条结论有两个直接用途：一是解释了为什么“延长周期省电”必须和瞬时峰值一起考虑（§7.3）；二是由于现行标准按小时均值考核，本文把瞬时峰值直接卡到限值的做法偏\textbf{保守}（§7.4）。中心情景取$\rho=1$，即假设四个电场同时振打、峰值完全叠加，这是最不利的上界；若错开相位则取$\rho=0.80$或0.65。敏感性分析仍扫描$\gamma\in\{1.00,1.30,1.60\}$，以检验推导值附近是否稳健。

\begin{figure}[htbp]
\centering
\includegraphics[width=0.92\textwidth]{figures_paper/14_rapping_peak_mechanism.png}
\caption*{\normalsize 图5　振打峰值机理与振打功率权衡}
\end{figure}

\subsection*{4.6 各因素方向性小结}

在上述骨架中，其他条件不变时各因素的作用方向如下。提高$U_i$会增大$K$、降低连续排放，但总功率按$U^2$上升。$K$不变时，入口浓度升高会使出口浓度同比升高。流量增大会缩短烟气在电场中的停留时间、同时抬高粉尘负荷，因而需要更强的捕集能力。温度则通过气体性质、粉尘比电阻和放电状态影响除尘效果，其中140～160 $^\circ$C的高比电阻区间最为不利（§4.3）。这些方向来自机理与文献[2-4]，不是从附件回归出来的。

\section{问题二：典型工况划分与10 mg/Nm³最低功率}

\subsection*{5.1 工况聚类}

聚类只用入口温度、入口浓度和烟气流量三个变量——它们由上游工艺决定，不受控制策略影响，因此不会出现“用结果去定义工况”的问题。三者标准化后作K-means聚类[9]。取$k=2,3,4,5,6$比较，轮廓系数[10]（衡量各类是否既紧凑又彼此分开，越大越好）依次为0.2963、0.3682、0.4084、0.3958、0.3934，在$k=4$最高。另一项Calinski–Harabasz指标依次为3204.2、4399.0、4721.2、4835.5、4712.0，最大值却出现在$k=5$，两项判据并不一致。

本文仍取$k=4$，理由有二：轮廓系数在$k=4$最高；四类中心恰好构成温度两档（约120.7与130.4 $^\circ$C）与流量两档（约44.0与48.1万 Nm³/h）的2×2组合，物理含义清楚。但必须说明，CH指标支持$k=5$，而0.4084的轮廓系数也只说明各类分得中等清楚。所以这四类是\textbf{统计上的分层，不是唯一正确的物理状态划分}；另外如§4.3所述，两段异常工况都落在四类中心之外。

\begin{figure}[htbp]
\centering
\includegraphics[width=0.92\textwidth]{figures_paper/03_cluster_selection.png}
\caption*{\normalsize 图6　聚类数选择结果}
\end{figure}

\begin{figure}[htbp]
\centering
\includegraphics[width=0.92\textwidth]{figures_paper/04_condition_map.png}
\caption*{\normalsize 图7　训练期典型工况分布}
\end{figure}


\begin{table}[htbp]
\centering
\caption*{\normalsize 表7　四类典型工况画像}
\vspace{-2pt}
\zihao{-5}
\begin{tabularx}{\textwidth}{>{\hsize=1.000\hsize\raggedright\arraybackslash}Xrrrrrrr}
\toprule
工况 & 占比 & 温度/$^\circ$C & 入口浓度/(g/Nm³) & 流量/(Nm³/h) & 粉尘负荷/(kg/h) & 历史总功率/kW & 锚点$K_0$ \\
\midrule
1 低温低浓度高流量 & 23.35\% & 120.65 & 26.08 & 478257 & 12460 & 1818.35 & 6.257 \\
2 高温中浓度低流量 & 20.06\% & 130.32 & 33.72 & 443591 & 14922 & 1870.89 & 6.514 \\
3 高温高浓度高流量 & 34.93\% & 130.41 & 40.71 & 480851 & 19560 & 1763.08 & 6.702 \\
4 低温高浓度低流量 & 21.67\% & 120.69 & 44.75 & 439990 & 19673 & 1696.55 & 6.797 \\
\bottomrule
\end{tabularx}
\end{table}

\subsection*{5.2 功率模型的设定与检验}

§4.1采用的形式并不是唯一可选的。为说明“选什么形式”本身有多重要，本文在同一时间划分下比较四种设定：只用均值的基线、电压一次项＋周期一次项的岭回归[11]、电压二次项＋周期一次项的岭回归，以及本文的$U_i^2+1/T_i$形式。岭回归的正则强度按训练期3个逐步加长的窗口滚动选择，本文形式则没有可调超参数。结果见表8。


\begin{table}[htbp]
\centering
\caption*{\normalsize 表8　功率模型设定比较与时间外推结果}
\vspace{-2pt}
\zihao{-5}
\begin{tabularx}{\textwidth}{>{\hsize=1.000\hsize\raggedright\arraybackslash}Xrrrrr}
\toprule
设定 & 验证$R^2$ & 验证RMSE/kW & 测试$R^2$ & 测试RMSE/kW & 测试平均偏差/kW \\
\midrule
均值基线 & −5.3259 & 140.18 & −0.3733 & 83.51 & −43.54 \\
电压一次＋周期一次（岭） & 0.9798 & 7.93 & 0.9180 & 20.41 & −16.65 \\
电压二次＋周期一次（岭） & 0.9837 & 7.12 & 0.9571 & 14.76 & −10.78 \\
\textbf{本文 $U_i^2+1/T_i$} & \textbf{0.9885} & \textbf{5.97} & \textbf{0.9930} & \textbf{5.98} & \textbf{−0.02} \\
本文设定＋工况变量 & 0.9885 & 5.97 & 0.9929 & 5.98 & +0.08 \\
\bottomrule
\end{tabularx}
\end{table}

这张表有三点值得强调。

第一，\textbf{把振打周期写成线性项是形式选错了，不是数据本身在漂移}。线性$T$的设定在测试日留下−10.78 kW的平均偏差，而且偏差在高功率区还会扩大；改成$1/T$后，偏差降到−0.02 kW，测试RMSE从14.76 kW降到5.98 kW，$R^2$从0.957升到0.993（图9）。如果沿用线性形式，这个偏差很容易被误当成“数据有问题、需要额外留误差预算”，实际上它是模型自己造出来的。

第二，$T$是优化时要调的量。线性形式会让模型以为“周期每延长1秒省下的电是固定的”，而真实的边际节电按$c_i/T_i^2$衰减——周期越长，再延长越不划算。用错形式会系统性地把最优周期算偏。

第三，再加入工况变量后各项指标不再改善，印证了§4.1的判断：总功率只由控制量决定。

验证集上绝对误差的90\%分位数为9.78 kW，本文把它作为做功率预算时的余量。它不改变各方案的优劣排序，也不保证在今后的日期里仍能覆盖90\%的情形。

\begin{figure}[htbp]
\centering
\includegraphics[width=0.92\textwidth]{figures_paper/07_power_prediction.png}
\caption*{\normalsize 图8　功率模型预测结果}
\end{figure}

\begin{figure}[htbp]
\centering
\includegraphics[width=0.92\textwidth]{figures_paper/08b_feature_importance.png}
\caption*{\normalsize 图9　振打周期的设定形式对比}
\end{figure}

\begin{figure}[htbp]
\centering
\includegraphics[width=0.92\textwidth]{figures_paper/08_power_residuals.png}
\caption*{\normalsize 图10　回顾性测试残差诊断}
\end{figure}

\begin{figure}[htbp]
\centering
\includegraphics[width=0.92\textwidth]{figures_paper/09_voltage_power_response.png}
\caption*{\normalsize 图11　局部电压—功率响应}
\end{figure}

\subsection*{5.3 优化模型}

设控制向量$z=(U_1,\ldots,U_4,T_1,\ldots,T_4)$。对每类工况$k$与限值$C_{lim}$求解

\begin{equation}
\min_z\ \widehat P(z)=a_0+\sum_i b_iU_i^2+\sum_i \frac{c_i}{T_i}
\end{equation}

\begin{equation}
\text{s.t.}\quad C_{base,k}(z)+B_k(z)\le C_{lim},\quad
U_1\ge U_3+3,\ U_2\ge U_4+3,\ T_3\ge T_1+100,\ T_4\ge T_2+100,\quad z\in\Omega_k 
\end{equation}

取值范围$\Omega_k$取该工况训练期内各控制量实际出现过的最小值到最大值。5 mg/Nm³一档另外允许电压上界放宽到历史最大值的1.05倍，放宽了多少在表11中逐条列出。已有的多电场节能研究同样采用“排放受限、分场优化”的思路[5,7]，但它们的排放模型都由现场实测标定，条件与本文不同，方法对照见附录A。

\textbf{关于求解方法}，有必要说明为什么不用现成的凸优化求解器。以电压$U$为变量时，目标函数$\sum b_iU_i^2$本身是凸的，但排放约束要求这样一个凸函数取值\textbf{足够大}——这类约束把可行域挖成了非凸的形状，整体不是凸规划。若改用$v_i=U_i^2$作变量，目标和排放约束都变成线性，可惜排序约束$U_1\ge U_3+3$在$v$下不再是凸集。也就是说，换哪种坐标都不能把它化成凸问题，宣称凸性会是错误的。

本文因此采用\textbf{多起点SLSQP确定性求解}：每个“工况—限值”组合从53个起点分别求解（5个有结构的起点：历史中位、上界、下界、中点、$T^*$点；再加48个随机起点），看它们是否收敛到同一个解。此外还在同一可行域内均匀撒$4\times10^5$个点做随机搜索，作为独立对照（§8.1）。这一做法取代了早期版本的百万级随机采样：问题结构清楚之后，确定性求解既更快，也更容易复算。

\subsection*{5.4 电场优先级的解析判据}

在给出数值解之前，先从模型结构直接推出“该先调哪个电场”。把$K_k=K_{0,k}\sum_i\alpha_i(U_i/U^{ref}_{i,k})^2/\sum\alpha$与$P=\sum_i b_iU_i^2$同时对$U_i^2$求导，得到\textbf{每多花一份电功率能换来多少去除指数}

\begin{equation}
\left.\frac{\mathrm{d}K}{\mathrm{d}P}\right|_i=\frac{K_{0,k}\,\alpha_i}{\left(\sum_j\alpha_j\right) b_i \left(U_{i,k}^{ref}\right)^{2}} 
\end{equation}

这个比值越大，说明该电场越“划算”。由于四个电场的$b_i$几乎相同，排序完全取决于$\alpha_i/(U^{ref}_{i,k})^{2}$，而\textbf{不是}取决于$\alpha_i$。道理是：功率按电压平方增长，参考电压本来就低的电场，抬高同样的比例所付出的电功率要少得多。代入中心权重与各工况的历史中位电压，得到表9。可以看到，第三电场的去除权重$\alpha_3=0.95$虽然低于前两级，但它的参考电压低约10 kV，平方之后使它的单位功率收益在四类工况中\textbf{都排第一}。


\begin{table}[htbp]
\centering
\caption*{\normalsize 表9　电场优先级解析判据（$\times10^{-3}$ kW$^-$$^1$）}
\vspace{-2pt}
\zihao{-5}
\begin{tabular}{lrrrrl}
\toprule
工况 & 电场1 & 电场2 & 电场3 & 电场4 & 优先次序 \\
\midrule
1 低温低浓度高流量 & 5.768 & 5.504 & \textbf{6.668} & 5.573 & E3＞E1＞E4＞E2 \\
2 高温中浓度低流量 & 5.556 & 5.269 & \textbf{7.023} & 5.870 & E3＞E4＞E1＞E2 \\
3 高温高浓度高流量 & 6.630 & 6.327 & \textbf{8.024} & 6.706 & E3＞E4＞E1＞E2 \\
4 低温高浓度低流量 & 7.551 & 7.207 & \textbf{7.717} & 6.502 & E3＞E1＞E2＞E4 \\
\bottomrule
\end{tabular}
\end{table}

这也解释了本文早期版本“前级优先”这一结论的真实来源：\textbf{它是排序约束$U_1\ge U_3+3$带出来的，而不是权重$\alpha_i$决定的}。因为抬高$U_3$就必须同时抬高$U_1$，两者被这条约束绑在一起，所以最优解里$U_1$和$U_3$总是同步走高。若把这条约束拿掉，最划算的仍然是第三电场。

振打侧同理：$\partial P/\partial T_i=-c_i/T_i^2$。前级周期短（中位240 s），后级周期长（中位441 s），所以\textbf{同样延长1秒，前级省下的电约为后级的3.4倍}。这与电压侧“前级承担主要捕集”并不矛盾：电压负责提高捕集能力，周期负责节省电耗，两者各自在不同的电场上最见效。

\subsection*{5.5 10 mg/Nm³下的最优控制参数}

四类工况的中心情景解如表10（5 mg/Nm³行为第四问提供参数）。


\begin{table}[htbp]
\centering
\caption*{\normalsize 表10　四工况两档目标的最优控制参数}
\vspace{-2pt}
\zihao{-5}
\resizebox{\textwidth}{!}{%
\begin{tabular}{lrrrrrrrrrrr}
\toprule
工况 & 目标/(mg/Nm³) & $U_1$/kV & $U_2$/kV & $U_3$/kV & $U_4$/kV & $T_1$/s & $T_2$/s & $T_3$/s & $T_4$/s & 峰值排放/(mg/Nm³) & 预测总功率/kW \\
\midrule
1 & 10 & 70.2 & 68.6 & 57.6 & 57.4 & 290 & 289 & 492 & 493 & 9.948 & 2045.55 \\
2 & 10 & 70.9 & 69.8 & 59.2 & 58.2 & 280 & 278 & 494 & 498 & 9.962 & 2104.30 \\
3 & 10 & 68.0 & 60.3 & 56.3 & 57.3 & 275 & 280 & 506 & 510 & 9.971 & 1924.67 \\
4 & 10 & 70.5 & 61.8 & 57.2 & 42.9 & 285 & 283 & 495 & 498 & 9.991 & 1851.15 \\
1 & 5 & 73.7 & 71.6 & 60.5 & 60.3 & 290 & 289 & 492 & 493 & 4.993 & 2182.25 \\
2 & 5 & 74.4 & 72.9 & 62.2 & 61.1 & 280 & 278 & 494 & 498 & 4.998 & 2245.51 \\
3 & 5 & 71.4 & 63.4 & 59.1 & 60.4 & 275 & 280 & 506 & 510 & 4.991 & 2056.06 \\
4 & 5 & 74.0 & 66.4 & 60.1 & 42.9 & 285 & 283 & 495 & 498 & 4.988 & 1973.10 \\
\bottomrule
\end{tabular}
}
\end{table}

\textbf{第二问的答案}：在10 mg/Nm³下，四类工况的最低总功率为1851.15～2104.30 kW，按训练期各工况占比加权为\textbf{1973.0 kW}，折合日电耗47.35 MWh。作为对照，历史实际加权功率是1783.2 kW（日电耗42.80 MWh）。也就是说，达到10 mg/Nm³意味着\textbf{功率上升10.6\%、日电耗多用4.55 MWh}。

这个方向是必然的：历史排放约50 mg/Nm³，要把出口浓度压到五分之一，就必须提高捕集能力；而在线可调的手段只有电压，电压又按平方计入功率。所以\textbf{“节能”只能在同一排放目标下、相对其他达标方案而言，不可能相对历史运行而言。}

解的结构与表9的判据一致。振打周期$T$在两档下都被推到历史区间的上端，因为延长周期直接降低振打功率；电压则按性价比高低分配——工况3与工况4中$U_2$明显低于$U_1$，工况4的$U_4$甚至被压到历史下界42.9 kV，正因为这个电场在该工况下最不划算（表9末行）。可见“四个电场一起抬”并不是最优做法。

表中数值是点预测值，做预算时可另加9.78 kW的误差余量。电压不必按0.1 kV的精度执行：多起点求解得到的功率完全一致，但控制量本身存在若干功率近乎相同的组合（§8.1）。

\subsection*{5.6 目标点已越出历史运行经验}

上表的解能否直接下发给设备？答案是明确的“不能”，而且理由可以用数据量化。

需要一把尺子来衡量“某个控制组合离历史常规操作有多远”。单看某一个电压是否超限并不够，因为八个控制量之间彼此相关——历史上从没有过“$U_1$很高而$U_3$很低”的组合，即便两者各自都在范围内。马氏距离平方$d_M^2=(z-\mu_k)^T\Sigma_k^{-1}(z-\mu_k)$正是把这种相关性一并计入的距离，其中$\mu_k$和$\Sigma_k$是该工况训练期控制量的均值与协方差。把它与历史上97.5\%的时间都不超过的水平$q_{k,0.975}$（四类工况依次为20.55、18.72、16.85、19.09）相比，就能看出推荐解偏离历史经验多远，结果见表11。


\begin{table}[htbp]
\centering
\caption*{\normalsize 表11　推荐解相对历史运行经验的外推审计}
\vspace{-2pt}
\zihao{-5}
\begin{tabular}{lrrrrc}
\toprule
工况与限值 & $d_M^2$ & $q_{0.975}$ & 倍数 & 最高电压超出历史最大值 & 判定 \\
\midrule
工况1，10 mg & 49.39 & 20.55 & 2.40× & 0.0\% & 外推 \\
工况2，10 mg & 46.46 & 18.72 & 2.48× & 0.0\% & 外推 \\
工况3，10 mg & 77.69 & 16.85 & 4.61× & 0.0\% & 外推 \\
工况4，10 mg & 102.79 & 19.09 & 5.38× & 0.0\% & 外推 \\
工况1，5 mg & 65.84 & 20.55 & 3.20× & +5.05\% & 外推 \\
工况2，5 mg & 60.31 & 18.72 & 3.22× & +5.07\% & 外推 \\
工况3，5 mg & 106.26 & 16.85 & 6.31× & +5.00\% & 外推 \\
工况4，5 mg & 133.43 & 19.09 & 6.99× & +5.07\% & 外推 \\
\bottomrule
\end{tabular}
\end{table}

结论有三层。

其一，最低功率解全都落在历史经验之外，$d_M^2$达到该水平的2.40～6.99倍。因此本文把$d_M^2$\textbf{作为外推风险指标报告}，而不把它当作必须满足的约束——否则求出来的就不是“最低功率”，而是“最低功率且最像以前的操作”，这是两个不同的问题。

其二，如果把$d_M^2\le q_{k,0.975}$真的\textbf{加成硬约束}，八个“工况—限值”组合里仍有七个有解（\texttt{results/mahalanobis\_constrained\_feasibility.csv}），代价是功率上升：10 mg/Nm³档四类工况分别多耗1.46\%～4.38\%（加权2026.5 kW，比最低功率解高2.7\%），5 mg/Nm³档三个有解的工况多耗2.22\%～6.07\%，只有工况1在5 mg/Nm³下无解。这就给出了一个可选的\textbf{保守方案}：如果现场不接受偏离既往操作习惯，可以多花几个百分点的电，换取“每一步都落在历史做过的组合里”。本文以最低功率解为主，同时保留这个保守方案，供现场按自身风险偏好选用。

其三，从电源容量看，10 mg/Nm³的四个解\textbf{仍在历史电压范围之内}——每个电场的电压都不超过该工况历史上出现过的最大值，说明这个目标在现有高压电源容量下原则上可以达到。而5 mg/Nm³需要各电场电压比历史最大值再高，最多高出5.07\%；四类工况所需的最小裕度分别是5\%、4\%、1\%、0\%。\textbf{这说明5 mg/Nm³首先是容量问题，其次才是电耗问题}：如果现场高压电源没有这点裕度，再怎么调度也达不了标，只能靠增设电场、改造电源或在前端减少入口粉尘来解决。

\begin{figure}[htbp]
\centering
\includegraphics[width=0.92\textwidth]{figures_paper/05b_boundary_proximity.png}
\caption*{\normalsize 图12　推荐解相对历史运行经验的外推审计}
\end{figure}

\subsection*{5.7 排放—功率权衡前沿}

把排放限值由10逐档收紧到5 mg/Nm³并重复求解，得到“排放—功率”权衡曲线（图13）。四类工况的最优功率随限值收紧单调上升，工况加权功率由1973.0 kW（10 mg/Nm³）升至2105.5 kW（5 mg/Nm³），中间各档依次为1991.3、2012.7、2046.3、2067.7 kW；所需电压裕度同步由0\%增至5\%（每收紧1 mg/Nm³约增1个百分点）。全部档位均高于历史实际加权功率1783.2 kW。这张图把“限值与功率”的关系整体呈现出来：\textbf{沿曲线向左走（限值更严）必然多耗电，向右走（限值更松）才省电}；而历史运行点还在这张图右侧很远的地方。

\begin{figure}[htbp]
\centering
\includegraphics[width=0.92\textwidth]{figures_paper/19_emission_power_frontier.png}
\caption*{\normalsize 图13　排放—功率权衡前沿}
\end{figure}

\section{问题三：代表工况的控制策略比较}

选取工况1和工况3。二者流量都较高，但入口浓度和粉尘负荷差异明显：工况1平均负荷12460 kg/h，工况3为19560 kg/h，可用于比较低负荷与高负荷策略。

\textbf{低负荷的工况1采取“四场一起抬”}：10 mg/Nm³下$U=(70.2,68.6,57.6,57.4)$ kV，四个电场相对各自参考值的抬升幅度接近。原因是这类工况的起算点$K_0=6.257$最低，要达标所需的提升倍数$S=1.288$最大，四个电场都得出力，没有挑选的余地。

\textbf{高负荷的工况3采取“挑着抬”}：$U=(68.0,60.3,56.3,57.3)$ kV，其中$U_1$、$U_3$、$U_4$都接近历史最大值，唯独$U_2$停在60.3 kV（它的历史最大值是69.3 kV）。这正是表9判据的直接体现——工况3的优先次序是E3＞E4＞E1＞E2，第二电场最不划算，因此被留在低位。限值由10收紧到5 mg/Nm³时，各场增量为$\Delta U=(+3.4,+3.1,+2.8,+3.1)$ kV，这一格局保持不变。

因此第三问的答案不是“前级优先”这样一句口号，而是三条可以直接算的规则：

\begin{itemize}[itemsep=1pt,topsep=3pt]
\item \textbf{电压侧}：按$\alpha_i/(b_iU^{ref\,2}_i)$降序抬升，受排序约束$U_1\ge U_3+3$牵引时同步抬升$U_1$；
\item \textbf{振打侧}：按$c_i/T_i^2$降序延长周期，即优先延长周期较短的前级；
\item \textbf{两者分工}：电压负责提高捕集能力，周期负责节省电耗，各自在不同电场上最见效，不应混为一谈。
\end{itemize}

历史数据对这套规则只提供了部分支持。§4.2的规律一（浓度高就缩短各场周期）在全部时段方向一致，与“高负荷下不宜为省电而延长周期”相符；规律二（浓度高就抬前两级电压）只在最后两天成立，因此本文只把它列作参照。

为便于现场执行，按排放构成与容量状态把上述比较整理成一张决策表，见表12。


\begin{table}[htbp]
\centering
\caption*{\normalsize 表12　控制优先级决策规则}
\vspace{-2pt}
\zihao{-5}
\begin{tabularx}{\textwidth}{>{\hsize=0.940\hsize\raggedright\arraybackslash}X>{\hsize=0.940\hsize\raggedright\arraybackslash}X>{\hsize=1.120\hsize\raggedright\arraybackslash}X}
\toprule
状态特征 & 优先动作 & 原因 \\
\midrule
连续排放基值接近限值，振打峰值占比较低 & 按$\alpha_i/(b_iU^{ref\,2}_i)$降序小步提高电压 & 主要矛盾是连续捕集强度不足，应优先动用性价比最高的电场 \\
峰值增量占比较高，周期明显长于参考周期 & 优先缩短过长周期并实施错相振打 & 主要矛盾是单次积灰脱落和峰值叠加 \\
周期过短、振打频繁，且电耗压力大 & 优先延长前级周期 & 前级$c_i/T_i^2$最大，单位延长的节电最多 \\
高粉尘负荷且基值和峰值均高 & 先错相、再调整周期，最后小步提升电压 & 先控制峰值，再补充连续捕集能力 \\
某电场电压已达历史最大值 & 转向次优电场；若各场均达上限则停止优化 & 数学可行不等于电源容量可行 \\
入口温度进入140～160 $^\circ$C高比电阻带 & 退出优化，按现场反电晕处置预案运行 & 该带未被本文工况覆盖，模型外推无效 \\
超出历史支持域或出现火花、联锁风险 & 停止优化，回退现场安全设定 & 数学可行不等于设备安全 \\
\bottomrule
\end{tabularx}
\end{table}

\begin{figure}[htbp]
\centering
\includegraphics[width=0.92\textwidth]{figures_paper/11_optimal_voltage.png}
\caption*{\normalsize 图14　两档排放限值下的电压配置}
\end{figure}

\begin{figure}[htbp]
\centering
\includegraphics[width=0.92\textwidth]{figures_paper/12_optimal_rapping_period.png}
\caption*{\normalsize 图15　两档排放限值下的振打周期}
\end{figure}

在线运行设\textbf{TARGET-10、TARGET-5、FALLBACK三种状态}。选哪一档由上级指令或管理规定决定，不设自动切换阈值——本文的排放模型未经现场标定，还不足以在线自动判定当前排放水平。TARGET-10与TARGET-5分别读取表10中对应工况、对应目标的控制参数。控制量按斜坡逐步逼近目标值，试运行阶段暂定电压变化不超过0.5 kV/min、振打周期不超过5 s/min；若现场DCS、联锁或设备制造商规定了更严的限制，以更严者为准。四个电场的振打先按$\rho=1$（同时振打、峰值完全叠加）这一最不利情形核验，再用$\rho=0.80$和0.65评估错开相位能带来多少余量。一旦传感器失效、火花率越限、工况超出取值范围或联锁触发，即切换到FALLBACK。

\section{问题四：10收紧至5 mg/Nm³的功率代价}

\subsection*{7.1 闭式解}

第四问可以直接解出公式，而公式比数值结果更有解释力。设在10 mg/Nm³的最优点上，电场功率为$P_{field}$、去除指数为$K_{10}$。把四个电场的电压同时放大$f$倍，则去除指数变成$K_{10}f^{p}$，电场功率变成$P_{field}f^{2}$。要让排放减半（计入峰值项后，所需的去除指数增量是$\Delta K=\ln\frac{10-B}{5-B}$），解得$f=(1+\Delta K/K_{10})^{1/p}$，于是

\begin{equation}
\Delta P=P_{field}\left[\left(1+\frac{\Delta K}{K_{10}}\right)^{2/p}-1\right]
\ \approx\ \frac{2}{p}\cdot\frac{P_{field}}{K_{10}}\,\Delta K 
\end{equation}

这个式子的含义很清楚：\textbf{限值每减半所增加的电耗，等于电场功率乘以所需的去除指数增量、除以当前的去除指数，再乘以$2/p$}。由此可见，设备当前越干净（$K$越大），再减半一次的相对代价反而越低。

\subsection*{7.2 分工况结果}


\begin{table}[htbp]
\centering
\caption*{\normalsize 表13　中心情景下10收紧至5 mg/Nm³的分工况功率代价}
\vspace{-2pt}
\zihao{-5}
\begin{tabularx}{\textwidth}{>{\hsize=1.000\hsize\raggedright\arraybackslash}Xrrrrrrr}
\toprule
工况 & 10 mg功率/kW & 5 mg功率/kW & $K_{10}$ & $\Delta K$ & 闭式$\Delta P$/kW & 数值$\Delta P$/kW & 功率增幅 \\
\midrule
1 低温低浓度高流量 & 2045.55 & 2182.25 & 8.059 & 0.804 & 137.7 & 136.7 & 6.68\% \\
2 高温中浓度低流量 & 2104.30 & 2245.51 & 8.338 & 0.832 & 142.4 & 141.2 & 6.71\% \\
3 高温高浓度高流量 & 1924.67 & 2056.06 & 8.594 & 0.903 & 131.3 & 131.4 & 6.83\% \\
4 低温高浓度低流量 & 1851.15 & 1973.10 & 8.690 & 0.906 & 123.1 & 122.0 & 6.59\% \\
\bottomrule
\end{tabularx}
\end{table}

闭式解与数值解的相对差异不超过1\%，互为验证。按训练期工况占比加权，

\begin{equation}
\bar P_{10}=1973.0\ \text{kW},\qquad \bar P_{5}=2105.5\ \text{kW},\qquad
\Delta P/\bar P_{10}=6.72\% 
\end{equation}

按相同运行时长折算，两档日电耗分别约为47.35与50.53 MWh/d，收紧限值的日增电耗约3.18 MWh/d；按年运行8760 h计，年增电耗约1.16×10³ MWh。加上9.78 kW验证期误差附加量后，两档加权功率分别为1982.8与2115.3 kW；由于两者使用同一附加量，该比例不作为严格上界，只用于功率预算。

值得注意的是，四类工况的增幅高度一致（6.59\%～6.83\%）。这不是巧合，而是上述公式的直接推论：增幅$\approx(2/p)(P_{field}/P)(\Delta K/K)$，而四类工况的电场功率占比和当前去除指数都相差不大，所以增幅几乎不随工况变化。

\begin{figure}[htbp]
\centering
\includegraphics[width=0.92\textwidth]{figures_paper/18_condition_energy_penalty.png}
\caption*{\normalsize 图16　分工况功率增幅比较}
\end{figure}

\subsection*{7.3 真正的权衡：振打周期能买回多少}

把第四问的答案与§4.1的振打结构放在一起看，可以得到本文最有操作价值的一组对照。按工况加权计算，限值由10收紧到5 mg/Nm³要多耗约132.5 kW；而把四个电场的振打周期从历史中位延长到历史P95可省下56.1 kW，延长到历史最大值可省下109.1 kW。\textbf{也就是说，振打这一侧能挖出的全部潜力（约109 kW），恰好抵掉排放收紧一档所增电耗（约133 kW）的82\%。}

但这笔账不能只算一半。按§4.5的推导，把周期从中位延长到历史最大值（平均延长17.9\%），单次振打的灰饼质量和瞬时峰值也会同比放大约17.9\%。而在表10的解中，峰值项本就已经占掉限值的0.95～1.61 mg/Nm³，相当于5 mg/Nm³这一档约五分之一到三分之一的余量。峰值再放大17.9\%，还要多吃掉0.17～0.29 mg/Nm³，这部分必须靠提高电压来补，而提高电压又要多花电。这就是“不能为省电而无限延长振打周期”的定量说法：\textbf{振打周期是这台设备上最便宜的节电手段，但额度有限，而且是用瞬时排放峰值换来的。}

\subsection*{7.4 峰值约束口径的说明}

在中心情景的解中，10 mg/Nm³方案的连续排放基础值为8.38～9.00 mg/Nm³，振打带来的峰值增量为0.95～1.61 mg/Nm³；5 mg/Nm³方案的基础值则降到3.38～4.05 mg/Nm³。可见只让平均值或基础值低于限值是不够的，否则振打的那一刻仍可能超标。

需要说明的是，本文把\textbf{瞬时峰值}直接卡到限值，而现行排放标准对颗粒物通常按小时均值考核。按§4.5的推导，一次持续数秒到数十秒的再飞扬峰值，对小时均值的贡献正比于$L$而与周期无关。因此本文的口径系统性偏保守，会高估达标所需的功率。若改用“小时均值达标、峰值只作软性惩罚”的口径，算出的功率增幅会低于本文报告值。在缺少逐次振打数据的情况下，本文宁可取保守口径，但不应把它当作唯一的合规解释。

\begin{figure}[htbp]
\centering
\includegraphics[width=0.92\textwidth]{figures_paper/13_emission_decomposition.png}
\caption*{\normalsize 图17　排放约束分解}
\end{figure}

\section{模型检验、敏感性与稳健性}

\subsection*{8.1 求解质量}

八个“工况—限值”组合中，多起点SLSQP有47～53个起点成功收敛，\textbf{所有收敛起点给出的最优功率一致到$10^{-7}$ kW以内}——说明它们都落到了同一个解上。作为独立对照，在同一可行域内均匀撒$4\times10^5$个点的随机搜索，在全部八个组合上都没有找到更好的解，其最优值反而高出0.97\%～5.17\%；其中工况1的10 mg/Nm³组合甚至一个可行样本都没采到。原因是可行域被排放约束和排序约束夹成了一条窄带，随机撒点几乎打不中。

这组对照说明：问题结构清楚之后，确定性求解优于大规模随机搜索。但仍要说明，对非凸问题而言，多起点收敛到同一点是\textbf{很强的证据，却不是数学证明}。因此本文只说“最优解”，不宣称“全局最优”。

\begin{figure}[htbp]
\centering
\includegraphics[width=0.92\textwidth]{figures_paper/15_search_convergence.png}
\caption*{\normalsize 图18　求解方法对照}
\end{figure}

\subsection*{8.2 先验敏感性}

功率增幅结论对灰箱先验的依赖集中在\textbf{单一标量$p$}上。由§7.1的闭式解，增幅严格正比于$2/p$；表14给出参考表。


\begin{table}[htbp]
\centering
\caption*{\normalsize 表14　功率增幅对幂次$p$的依赖}
\vspace{-2pt}
\zihao{-5}
\begin{tabularx}{\textwidth}{l>{\hsize=1.000\hsize\raggedright\arraybackslash}Xrc}
\toprule
$p$ & 物理含义 & 10$\rightarrow$5 mg/Nm³加权增幅 & 5 mg/Nm³在5\%电压裕度内是否可达 \\
\midrule
1.0 & 去除指数正比于电压 & 14.2\% & 否，需更大电压裕度 \\
1.5 & 中间情形 & 9.2\% & 否，需更大电压裕度 \\
\textbf{2.0} & \textbf{Deutsch迁移速度形式（中心）} & \textbf{6.7\%} & \textbf{是} \\
2.5 & 强电压响应 & 5.3\% & 是 \\
3.0 & 极强电压响应 & 4.4\% & 是 \\
\bottomrule
\end{tabularx}
\end{table}

\textbf{$p$是本文结论中唯一需要现场标定的数，而且一次电压阶跃试验就能测出来}：在工况稳定时把四个电场的电压同时抬高几个百分点，记录未截顶的出口浓度如何变化，再把$\ln C_{out}$对$\ln U$作回归，斜率就是$-K_0 p$。对现场来说，给出这样一个可测的量，比丢下一个很宽的区间有用得多。

表中增幅为闭式解值；对$p\ge2$的三档另作数值重解校验（图20叉号），与闭式值相差不超过0.4个百分点。

除$p$以外，本文还扫描了两组情景。

一是取值大小的组合：$p\in\{1.5,2.0,2.5\}$×峰值尺度$b_0\in\{0.9,1.2,1.5\}$×安全裕度$s\in\{0,0.08,0.15\}$，共27组。其中15组在5\%电压裕度内四类工况全部有解；无解的12组全部来自$p=1.5$，这再次说明$p$首先决定的是\textbf{能不能达标}，而不只是代价大小。有解各组的加权增幅为3.99\%～7.04\%，中位数5.85\%。

二是模型形状的组合：电场权重形状（3种）×峰值指数$\gamma\in\{1.00,1.30,1.60\}$×相位尺度$\rho\in\{0.65,0.80,1.00\}$×参考周期负荷指数（3种），共81组，全部有解，加权增幅为5.13\%～7.55\%，中位数6.40\%。峰值指数在推导值1.00附近变化时增幅几乎不变，这与§4.5的推导一致：峰值项占用的是限值余量，本身不是主导项。

需要强调，以上区间都是\textbf{在人为设定的参数网格上算出的情景范围，不是统计意义上的置信区间}。

\begin{figure}[htbp]
\centering
\includegraphics[width=0.92\textwidth]{figures_paper/16_sensitivity_distribution.png}
\caption*{\normalsize 图19　结构情景下的功率增幅分布}
\end{figure}

\begin{figure}[htbp]
\centering
\includegraphics[width=0.92\textwidth]{figures_paper/20_eta_sensitivity.png}
\caption*{\normalsize 图20　功率增幅对幂次$p$的依赖}
\end{figure}

\begin{figure}[htbp]
\centering
\includegraphics[width=0.92\textwidth]{figures_paper/17_sensitivity_heatmap.png}
\caption*{\normalsize 图21　电场贡献权重消融结果}
\end{figure}

\subsection*{8.3 稳健性结论}

稳健性证据分四层。第一，轮廓系数支持把工况分成四类。第二，四种设定的比较与时间外推支持所选的功率模型，其中$1/T$形式在\textbf{完全未参与拟合}的测试日上表现最好，这是本文最强的实证证据。第三，多起点收敛一致以及随机搜索的对照，支持数值解的质量。第四，闭式解与数值解互相验证，情景扫描与权重消融则检验了结论对建模选择的依赖程度。

但仍须强调：排放层的全部数值都是机理推演的结果。无论是误差余量还是情景范围，都不能把一个机理模型变成经过实测验证的模型。

\section{模型评价、实施方案与局限}

\subsection*{9.1 模型优点}

第一，先逐列判断数据能支持什么、不能支持什么，既避免用复杂算法去拟合截顶噪声，也没有因为一列不可用就放弃其余13列的信息。

第二，功率模型采用有物理含义的形式而非通用多项式，在未参与拟合的测试日上达到$R^2=0.993$、偏差−0.02 kW，并由此发现“振打占电耗近四成”这一可直接使用的结论。

第三，排放起算点直接取记录值，不引入未经证实的量纲假设，使第二问回答的是题目实际提出的问题。

第四，电场优先级由公式推出，而不是从某一次求解的现象里归纳。

第五，第四问给出闭式解，把全部不确定性收敛到一个可现场标定的数$p$。

第六，把“结果偏离历史运行经验多远”当作结论的一部分报告，并给出留在历史经验之内的保守方案及其电耗代价，而不是回避。

\subsection*{9.2 模型局限}

最主要的局限是，排放层的全部系数（$p,\alpha_i,\beta_i,b_0,\gamma,\rho,s$）都无法从附件中辨认，只能依据机理和文献设定；表14给出的区间是人为参数网格上的情景范围，不是统计置信区间。

其次，历史运行范围并不等于设备的安全边界，5 mg/Nm³档所需的5\%电压外推更须经高压电源容量与绝缘裕度核算。模型也没有二次电流、分场功率、火花率、极板积灰状态和联锁信息；分钟级数据不足以观察逐次振打峰值，因此相位叠加尺度$\rho$无法标定。四类工况的温度中心都没有覆盖140～160 $^\circ$C的高比电阻区间，本文结论在该区间无效。此外，测试日的数据曾被前期处理流程接触过，不能算作完全未触碰的盲测。

因此，本文给出的控制参数只能作为现场试验的起点，\textbf{不得直接下发到生产设备}。

\subsection*{9.3 上线与补采方案}

建议分四步部署。

\textbf{第一步：只算不动。} 系统在后台跟着算，但不改设定，用来验证工况识别是否正确、TARGET-10/TARGET-5参数表是否正确加载；同时完成高压电源容量核算，确认5 mg/Nm³所需的电压裕度是否真的存在。

\textbf{第二步：小步试验。} 在每类工况内做单因素小步阶跃，重点是\textbf{标定幂次$p$}（§8.2），并由现场确定电压与周期的真实变化速率上限。

\textbf{第三步：重新标定。} 用1～5 s的未截顶出口浓度、实际振打事件、二次电流、火花率和分场功率，重新标定排放层参数和相位尺度$\rho$，并单独补采140～160 $^\circ$C高比电阻区间的数据。

\textbf{第四步：盲测后闭环。} 冻结模型，在全新日期上盲测，通过后才允许投入闭环。任何超出取值范围、传感器失效或联锁事件都应触发FALLBACK。EPA技术资料与相关实验都强调气流分布、收尘面积和再飞扬对性能评价的重要性[2-3,6]，因此补采数据时应优先补齐这几项。

\section{结论}

\begin{enumerate}[itemsep=1pt,topsep=3pt]
\item \textbf{问题一：分两层回答。} 能由数据确定的部分：总功率几乎完全由控制量决定，$P=73.1+0.0850\sum U_i^2+5.39\times10^4\sum T_i^{-1}$（测试日$R^2=0.993$），其中振打环节在中位工况下占总功率的38.8\%，而工况变量对总功率没有额外解释力；历史操作规律中，“浓度高就缩短各场振打周期”在四个时段方向一致，“浓度高就抬前级电压”却只在最后两天成立；附件还含两段持续异常工况，其中05-06的158.2 $^\circ$C事件正落在粉尘比电阻的高值区。不能确定的部分：出口浓度记录在50 mg/Nm³附近被量程截顶，与全部12个变量的相关系数绝对值都小于0.017，无法从数据实证出“控制如何影响排放”。振打峰值则由积灰量守恒推出：单次峰值随周期近似线性放大，而它对小时均值的贡献与周期无关。
\item \textbf{问题二：4类工况与10 mg/Nm³方案。} 用入口三变量把运行状态分成4类（轮廓系数0.4084）。以记录值50 mg/Nm³为起算点，四类工况达到10 mg/Nm³的最低功率为1851.15～2104.30 kW，加权1973.0 kW，日电耗47.35 MWh，比历史实际的1783.2 kW高10.6\%。收紧排放必然多耗电，“排放不变还省电”并不存在。该方案的电压仍在历史范围之内，但离历史常规操作已相当远（马氏距离达97.5\%水平的2.4～5.4倍）；若要求必须留在历史经验之内，四类工况仍都有解，加权功率升到2026.5 kW，多耗2.7\%。
\item \textbf{问题三：优先级可以直接算出来。} 每份电功率换来的除尘收益是$K_0\alpha_i/(\sum\alpha\cdot b_iU^{ref\,2}_i)$。第三电场的参考电压最低，因而性价比在四类工况中都排第一；“前级优先”的说法其实是排序约束$U_1\ge U_3+3$带出来的。振打侧按$c_i/T_i^2$排序，应优先延长周期较短的前级。低负荷工况必须四场一起抬，高负荷工况则可以挑着抬，让最不划算的电场留在低位。
\item \textbf{问题四：多花的电与真正的瓶颈。} 限值由10收紧到5 mg/Nm³时，闭式解$\Delta P=P_{field}\ln[(10-B)/(5-B)]/K_{10}$给出123.1～142.4 kW，与数值解相差不足1\%；加权功率由1973.0升到2105.5 kW，增幅6.72\%，日电耗由47.35增至50.53 MWh。增幅与$2/p$成正比，而$p$是全文唯一需要现场标定的数。更关键的限制不是电耗而是电源容量：5 mg/Nm³要求各电场电压比历史最大值再高，最多高出5.07\%，四类工况所需的最小裕度分别为5\%、4\%、1\%、0\%；若电源没有这点裕度，就只能靠增设电场、改造电源或在前端减少入口粉尘来解决，靠调度是解决不了的。以上排放数值均属机理推演，必须在补采未截顶的排放数据、振打事件与电气状态后重新标定。
\end{enumerate}

\section*{参考文献}
\addcontentsline{toc}{section}{参考文献}

[1] 中华人民共和国生态环境部. 水泥工业大气污染物排放标准：GB 4915—2013（含2025年修改单）[EB/OL]. (2013-12-27)[2026-08-06]. \url{https://www.mee.gov.cn/ywgz/fgbz/bz/bzwb/dqhjbh/dqgdwrywrwpfbz/201312/t20131227_265765.htm}.

[2] OGLESBY S Jr, NICHOLS G B. *A Manual of Electrostatic Precipitator Technology: Part I—Fundamentals*[R/OL]. Birmingham: Southern Research Institute, 1970[2026-08-06]. \url{https://nepis.epa.gov/Exe/ZyPURL.cgi?Dockey=9101FJFA.TXT}.

[3] SZABO M F, SHAH Y M, SCHLIESSER S P. *Inspection Manual for Evaluation of Electrostatic Precipitator Performance*[R/OL]. Washington, DC: U.S. Environmental Protection Agency, 1981[2026-08-06]. \url{https://nepis.epa.gov/Exe/ZyPURL.cgi?Dockey=9400034C.TXT}.

[4] MIZUNO A. Electrostatic precipitation[J]. IEEE Transactions on Dielectrics and Electrical Insulation, 2000, 7(5): 615-624. DOI: \url{10.1109/94.879357}.

[5] LI Z, SHAO C, AN Y, et al. Energy-saving optimal control for a factual electrostatic precipitator with multiple electric-field stages based on GA[J]. Journal of Process Control, 2013, 23(8): 1041-1051. DOI: \url{10.1016/j.jprocont.2013.06.007}.

[6] NAVARRETE B, VILCHES L F, RODRIGUEZ-GALAN M, et al. A pilot scale study of the rapping reentrainment and fouling in electrostatic precipitation[J]. Environmental Progress \& Sustainable Energy, 2015, 34(1): 7-14. DOI: \url{10.1002/ep.11934}.

[7] ZHENG C, ZHAO Z, GUO Y, et al. A real-time optimization method for economic and effective operation of electrostatic precipitators[J]. Journal of the Air \& Waste Management Association, 2020, 70(7): 708-720. DOI: \url{10.1080/10962247.2020.1767227}.

[8] FU Y, SU Z. Deep neural network based concentration prediction model of dry electrostatic precipitator with a modified PID optimizer[J]. Journal of Physics: Conference Series, 2022, 2189: 012018. DOI: \url{10.1088/1742-6596/2189/1/012018}.

[9] MACQUEEN J. Some methods for classification and analysis of multivariate observations[C]//Proceedings of the Fifth Berkeley Symposium on Mathematical Statistics and Probability. Berkeley: University of California Press, 1967, 1: 281-297.

[10] ROUSSEEUW P J. Silhouettes: A graphical aid to the interpretation and validation of cluster analysis[J]. Journal of Computational and Applied Mathematics, 1987, 20: 53-65. DOI: \url{10.1016/0377-0427(87)90125-7}.

[11] HOERL A E, KENNARD R W. Ridge regression: Biased estimation for nonorthogonal problems[J]. Technometrics, 1970, 12(1): 55-67. DOI: \url{10.1080/00401706.1970.10488634}.

[12] WHITE H J. *Industrial Electrostatic Precipitation*[M]. Reading, MA: Addison-Wesley, 1963.

[13] PARKER K R. *Electrical Operation of Electrostatic Precipitators*[M]. London: Institution of Electrical Engineers, 2003.

\section*{附录A 与既有多电场节能优化研究的对照}
\addcontentsline{toc}{section}{附录A 与既有多电场节能优化研究的对照}

本文与文献[5]、[7]同属“在排放约束下对多个电场分别优化”这一类工作，但两点不同：排放模型怎么标定，以及可行范围怎么划。对照见表A1。


\begin{table}[htbp]
\centering
\caption*{\normalsize 表A1　与既有多电场节能优化研究的方法对照}
\vspace{-2pt}
\zihao{-5}
\begin{tabularx}{\textwidth}{l>{\hsize=0.828\hsize\raggedright\arraybackslash}X>{\hsize=0.770\hsize\raggedright\arraybackslash}X>{\hsize=1.402\hsize\raggedright\arraybackslash}X}
\toprule
维度 & 文献[5]（GA优化） & 文献[7]（实时优化） & 本文 \\
\midrule
排放模型标定 & 由现场实测出口浓度回归 & 由在线仪表实时校正 & 出口浓度不可识别，改以记录值为锚点的Deutsch灰箱，全程标注为推演 \\
功率模型 & 分场电源功率实测 & 分场功率与二次电流 & 仅有总功率，按$U^2+1/T$物理设定辨识，测试$R^2=0.993$ \\
振打处理 & 未纳入优化 & 作为固定运行参数 & 纳入决策变量，并由积灰量守恒推导峰值标度 \\
求解方法 & 遗传算法 & 滚动实时优化 & 多起点SLSQP＋解析优先级判据＋随机搜索对照 \\
可行域 & 设备额定范围 & 在线安全约束 & 历史运行包络，越界量作为结论显式报告 \\
\bottomrule
\end{tabularx}
\end{table}

这三项工作的共同前提，是排放模型可以由现场数据标定。本文缺少这个条件，处理办法是把无法标定的部分全部收敛到一个幂次$p$上，并给出具体的标定方案（§8.2）——这是本文与上述工作的主要差别。

\section*{附录B 以出口浓度为直接目标的对照建模}
\addcontentsline{toc}{section}{附录B 以出口浓度为直接目标的对照建模}

为便于与其他方案交叉比较，这里给出直接以\texttt{C\_out\_mgNm3}为目标建模的结果（训练/验证仍按时间划分）：

\begin{enumerate}[itemsep=1pt,topsep=3pt]
\item 以“是否等于50”为标签做分类，验证集AUC=0.5098——与抛硬币无异；
\item 只用未截顶的样本做回归，验证集$R^2=-0.0190$；
\item 用全部记录值做回归，验证集$R^2=-0.0219$。
\end{enumerate}

后两项的$R^2$为负，意味着模型还不如直接取一个常数。因此，任何在这一列上报出高$R^2$的模型，其性能几乎必然来自两种情况之一：随机划分数据造成的时间泄漏，或者只是复现了那5491条恒为50的封顶值。

\section*{附录C 结果适用性声明}
\addcontentsline{toc}{section}{附录C 结果适用性声明}

本文的功率模型、工况划分、历史操作规律、异常工况识别以及数据质量统计，均直接来自附件，可用\texttt{src/}中的程序复算。

排放层的各项系数、振打峰值系数，以及由此得到的5/10 mg/Nm³控制参数和功率增幅，则属于机理推演。若数据提供方给出新的未截顶排放数据，或现场标定得到的幂次与$p=2$不同，应按§8.2重新计算，不得沿用本文的数值。5 mg/Nm³档所需的电压外推，须先经高压电源容量与绝缘裕度核算。

\section*{附录D 支撑材料与AI工具使用声明}
\addcontentsline{toc}{section}{附录D 支撑材料与AI工具使用声明}

\textbf{求解流程：}

\begin{lstlisting}[style=pseudo]
输入：分钟级数据、排放限值C_lim
1. 按时间划分训练/验证/回顾性测试；
2. 逐列可识别性审计；提取历史操作规律与异常工况；
3. 在训练期对Temp、C_in、Q标准化，K-means取k=4；
4. 比较四种功率设定，选定 P = a0 + sum b_i U_i^2 + sum c_i / T_i；
5. 以记录值50 mg/Nm³为锚点建立灰箱排放层；由积灰量守恒定 gamma=1；
6. 对每个工况与限值：
   6.1 取该工况控制量的历史最小—最大范围为支持域；
   6.2 多起点SLSQP求解（5结构起点＋48随机起点）；
   6.3 舍入至0.1 kV / 1 s 后按解析优先级修复可行性；
   6.4 以4×10^5点随机搜索独立对照；
7. 报告马氏距离外推审计与所需电压裕度；
8. 扫描27组尺度情景与81组结构情景，给出幂次p参考表；
输出：控制参数、功率代价、外推审计、敏感性、图表与运行摘要。
\end{lstlisting}

\textbf{源程序清单：} \texttt{src/scenario\_model.py}（功率模型、灰箱排放情景、支持域优化与敏感性）、\texttt{src/paper\_figures.py}（21幅正文图与图注追踪表）、\texttt{src/plot\_utils.py}（统一绘图样式）、\texttt{src/integrity\_check.py}（确定性核验）、\texttt{src/compute\_frontier.py}（排放—功率前沿）、\texttt{src/build\_paper\_pdf.py}（提交稿装配）。生成本文全部结果与图表的完整可运行源程序见附录E，与电子支撑材料 \texttt{src/} 目录逐字一致。

\textbf{支撑材料清单：} \texttt{data/}处理数据包；\texttt{results/}下功率基线、滚动验证、最优控制、外推审计、电压裕度需求、求解对照、结构情景、场权重消融、幂次参考表、动态控制表和运行摘要；\texttt{figures\_paper/}下21幅正文图（20幅PNG/PDF由绘图程序生成，图1为TikZ源文件独立编译）及\texttt{figure\_manifest.csv}；\texttt{integrity/}下确定性核验结果。

\textbf{AI工具使用声明（提交前由参赛队核验确认）：} 本稿坚持参赛队人工主导。题意分析、模型构建、公式推导、参数设定、程序审阅、结果复算和论文定稿由参赛队承担；AI编码助手仅用于代码排错、图表版式检查、文字校对与格式核验。AI输出均经人工筛选和复核，未替代现场数据，也未把不可识别的排放关系改写为实证结论。
\clearpage
\section*{附录E\quad 完整源程序}
\addcontentsline{toc}{section}{附录E\quad 完整源程序}
本附录给出生成本文全部结果与图表的完整可运行源程序，与电子支撑材料 \texttt{src/} 目录逐字一致。按依赖顺序排列：情景模型与优化、排放—功率前沿、绘图样式与图集生成、确定性核验、提交稿装配。
\subsection*{E.1\quad \texttt{scenario\_model.py}}
\lstinputlisting[style=pycode]{/home/user/Mathematical-Modeling/04_最终论文与代码交付包/src/scenario_model.py}
\subsection*{E.2\quad \texttt{paper\_figures.py}}
\lstinputlisting[style=pycode]{/home/user/Mathematical-Modeling/04_最终论文与代码交付包/src/paper_figures.py}
\subsection*{E.3\quad \texttt{plot\_utils.py}}
\lstinputlisting[style=pycode]{/home/user/Mathematical-Modeling/04_最终论文与代码交付包/src/plot_utils.py}
\subsection*{E.4\quad \texttt{compute\_frontier.py}}
\lstinputlisting[style=pycode]{/home/user/Mathematical-Modeling/04_最终论文与代码交付包/src/compute_frontier.py}
\subsection*{E.5\quad \texttt{integrity\_check.py}}
\lstinputlisting[style=pycode]{/home/user/Mathematical-Modeling/04_最终论文与代码交付包/src/integrity_check.py}
\subsection*{E.6\quad \texttt{build\_paper\_pdf.py}}
\lstinputlisting[style=pycode]{/home/user/Mathematical-Modeling/04_最终论文与代码交付包/src/build_paper_pdf.py}

\end{document}
